\cotefile{62}
\scanpage{68}
Expression au moyen de la fonction $\sigma$
\[
\boxed{\mathrm{sn}\,z=\sigma\omega_1\,e^{\eta_1 z}\frac{\sigma z}{\sigma(z+\omega_1)}
=-\sigma\omega_1\,e^{-\eta_1 z}\frac{\sigma z}{\sigma(z-\omega_1)}} \tag*{(6)}
\]

\unsure{On a|On|$q$} \deleted{$\mathrm{sn}\,z$ admet $\sigma(z$}

On en déduit
\[
\left\{\begin{aligned}
\mathrm{sn}(z+2\omega_1)&=\mathrm{sn}\,z\\
\mathrm{s}\deleted{\unsure{m|n}}\mathrm{n}(z+2\omega_2)&=-\mathrm{sn}\,z\\
\mathrm{sn}(z+2\omega_3)&=-\mathrm{sn}\,z\\
\text{d'où}\qquad\mathrm{sn}(z+4\omega_2)&=\mathrm{sn}\,z\\
\mathrm{sn}(z+4\omega_3)&=\mathrm{sn}\,z
\end{aligned}\right. \tag*{(7)}
\]
Donc $\mathrm{sn}\,z$ admet les périodes $2\omega_1$ et $4\omega_2$ (et $4\omega_3$)

Ses poles sont $z=\omega_1+2n_1\omega_1+4n_{\unsure{1|2}}\omega_2$, ses zéros $z=2n_1\omega_1+4n_1\omega_2$, poles et zéros sont simples ; $\mathrm{sn}\,z$ est une fonction d'ordre 2.

Ses zéros sont homologues de $0$ ou $2\omega_{2}$, où la dérivée est respectivement $1$ et $-1$ ; les poles sont homologues de pole $\omega_1$ et $\omega_1+2\omega_2$, dont les résidus sont resp\textsuperscript{t} $\frac{1}{k}$ et $-\frac{1}{k}$ (en supposant toujours choisi \inserted{le signe de} $k$ tel que l'on ait (5), ou, ce qui revient au même $\displaystyle k=-\frac{e^{\eta_1\omega_1}}{\sigma^{\unsure{2|1|\prime}}\omega_1}$)

On voit sur 3 ou 4 que $\mathrm{sn}\,z$ est paire ou impaire ; elle n'est pas \deleted{im}paire, car on aurait $\mathrm{sn}\,\omega_2=\mathrm{s}\deleted{\unsure{n}}\mathrm{n}\,{-\omega_2}$, mais aussi $\mathrm{sn}\,\omega_2=-\mathrm{sn}\,{-\omega_2}$, d'où $\mathrm{sn}\,\omega_2=0$, ce qui n'est pas. On peut aussi s'en assurer sur 6
\[
\boxed{\mathrm{sn}\,{-z}=-\mathrm{sn}\,z}
\]

Sur 3 ou 4, on voit $\mathrm{sn}^2\omega_2=\frac{1}{\deleted{\unsure{k^2|k}}}$ ; en choisissant convenablement le signe de $\omega_2$, on aura \uline{$\mathrm{sn}\,\omega_2=\frac{1}{\deleted{\unsure{k}}}$} ; l'équation différentielle vérifiée par $\mathrm{sn}\,z$ donne \uline{$\mathrm{sn}'\omega_2=0$} de $=$ $\mathrm{sn}\,{-\omega_2}=-1\quad\mathrm{sn}'{-\omega_2}=0$.

On voit alors sur 4 \deleted{que}, puis sur l'équation différentielle, que $\mathrm{sn}(\omega_1+\omega_2)=\frac1k$ soit \uline{$\mathrm{sn}\,\omega_3=-\frac1k\quad\mathrm{sn}'\omega_3=0$}.

% archive sheet number: 67
